\documentclass[11pt,a4paper]{article}
\usepackage[margin=27mm]{geometry}
\usepackage{amsmath,amsthm,amssymb,fontspec,unicode-math,hyperref}
\usepackage{longtable,booktabs,array,calc,fvextra}
\setmainfont{STIXGeneral}
\setmathfont{STIX Two Math}
\setmonofont{Menlo}
\hypersetup{hidelinks}
\newtheorem{theorem}{Theorem}
\newtheorem{proposition}[theorem]{Proposition}
\newtheorem{lemma}[theorem]{Lemma}
\newtheorem{corollary}[theorem]{Corollary}
\theoremstyle{remark}\newtheorem{remark}[theorem]{Remark}
\newcommand{\gen}[1]{\langle #1\rangle}
\newcommand{\F}{\mathbb F}
\newcommand{\Z}{\mathbb Z}
\newcommand{\Aut}{\operatorname{Aut}}
\newcommand{\Sym}{\operatorname{Sym}}
\setlength{\emergencystretch}{2em}
\providecommand{\tightlist}{\setlength{\itemsep}{0pt}\setlength{\parskip}{0pt}}

\title{Changes of chamber generators and binary translation groups}
\author{Human Compression Project}
\date{}
\begin{document}
\maketitle
\begin{abstract}
We study changes of distinguished chamber generators using subgroup
intersections, coset representatives and the objects retained by the
operation. Endpoint conjugation has an exact intersection criterion;
its geometric realization is a component of a graph's bipartite double
cover. We also give a common algebraic description of the binary translations that
occur in halving generalized cubes and in polyhedra with equal adjacent parameters.
Edge differences in a connected graph present its augmentation module;
an explicit covariance identity then identifies groups generated by
these translations with semidirect products. For the cyclic displacement
relations used by Li and Zhang, this gives mutually inverse homomorphisms
and proves the universal presentation without comparing finite group
orders. The arguments isolate the graph generation and presentation
recognition steps so that geometric applications can use the same
statements with their own distinguished generators.
\end{abstract}
\section{Introduction}
Many constructions by involutions have a translation subgroup whose
elements are even subsets of a finite set. The fact that the subset is
even explains its size, but it does not by itself prove that the proposed
generators produce every translation, or that an abstract presentation
has no additional kernel. These two questions have a common answer in
terms of differences of vertices in a connected graph.

We distinguish the underlying module from its chosen presentation.
Paths express every difference from a base vertex as a sum of edge
differences; cycle relations are exactly the ambiguity in this
expression. Once a group permutes the vertices, the corresponding
conjugation formula determines the whole semidirect product. This gives
an explicit inverse map in the cyclic application, rather than a
surjection followed by an equality of orders.

The binary-coordinate construction and its applications to generalized
cubes are attributed to Danzer and to Montero--Ivi\'c Weiss. The
commutativity of the cyclic displacements is an expanded argument from
Li--Zhang. The formulation below reorganizes these existing ingredients
into statements that can be used by both applications. No historical
priority is asserted for graph cycle spaces, augmentation modules or
semidirect products.
The geometric statements use two distinct hypotheses. Intersections of
subgroups generated by subsets of markers identify flag stabilizers;
a common point for pairwise incident cosets supplies flag completion.
The two hypotheses are kept separate in the recognition theorem.
For an operation retaining vertices or edges, its action kernel is
computed on those actual retained objects, rather than inferred from
an abstract presentation. This distinction is needed by Petrial
operations, whose natural action need not be faithful.

The graph-halving construction, the lattice interpretation of
nondegeneracy and the index-two criterion are attributed to the
halving literature, including [PT24]. The endpoint support criterion,
exact parity compatibility and retained-object formulation reuse the
project's previously reviewed complete arguments. We include their
proofs here: another topic's possession of them does not reduce this
topic's common-theory cost.
\section{Coset geometries and the objects retained by an
operation}\label{coset-geometries-and-the-objects-retained-by-an-operation}

A geometry has a finite set of types; its flags are sets of pairwise
incident objects of distinct types, every flag extends to a chamber, and
its residues are the induced geometries on objects incident with a fixed
flag. It is thin when every panel has exactly two completions and
residually connected when every residue of rank at least two has
connected incidence graph. Regularity means chamber transitivity of the
full type-preserving automorphism group. No finiteness of the object set
is implicit.

A marking of a group G in the geometric sections is a finite indexed generating family of nonidentity
involutions \((r_i)_{i\in I}\). Write \(G_J=\langle r_j:j\in J\rangle\).
The intersection property is \[
G_J\cap G_K=G_{J\cap K}\qquad(J,K\subseteq I).
\tag{C1}
\] The maximal parabolics are \(P_i=G_{I\setminus\{i\}}\); their right
cosets are incident when they intersect. Typed coset Helly means that
every pairwise intersecting family with distinct types has a common
element. UC.Geometry shows that IP and typed Helly give exactly the
regular thin residually connected coset geometry, with G as its full
type-preserving group. IP alone is not sufficient for that conclusion.

\textbf{Proposition UC.Geometry (coset recognition and faithful chamber
action).} Suppose the marking of \(G\) consists of nonidentity
involutions, has IP, and has typed coset Helly for its maximal
parabolics. The coset incidence system is then a thin, residually
connected geometry, and \(G\) is its full type-preserving automorphism
group, acting regularly on chambers. More generally, the full
type-preserving automorphism group of any thin, residually connected
geometry of finite rank acts freely on chambers. A chamber-transitive
subgroup is consequently the full group.

\emph{Proof.} In the coset system, a common coset point extends every
flag to a chamber by choosing the coset of each missing type through
that point. Right translation sends every such flag to its base flag.
For a base flag of type set \(S\), IP gives stabilizer
\(G_{I\setminus S}\). In particular the full chamber stabilizer is
\(1\), while a panel stabilizer is a group generated by one nonidentity
involution and has order two. There are therefore exactly two chambers
through each panel.

Consider a residue with remaining type set \(T\), where \(|T|\geq2\).
Its stabilizer is \(G_T\), and its maximal object stabilizers are
\(G_{T\setminus\{i\}}\), \(i\in T\). Every marker indexed by \(T\) lies
in at least one of these proper parabolics. A word in those subgroups
supplies a walk between their translated base objects in the residue
incidence graph; all objects are translates of base objects and all base
objects lie in a common flag. The residue is therefore connected.

For the general assertion, first observe that every residually connected
finite-rank geometry has a connected chamber graph. This follows by
induction on rank: chambers through a fixed object are connected through
galleries in its residue, and a path in the incidence graph, extending
each consecutive incident pair to a chamber, connects arbitrary chambers
through these subgraphs. The rank-one case holds because all chambers
share the empty panel. In a thin geometry an automorphism fixing a
chamber fixes its unique neighbor across each panel. It therefore fixes
every chamber along every gallery, and hence every object. The chamber
action is free. If a subgroup is chamber-transitive, any full
automorphism agrees on a base chamber with one subgroup element; their
quotient fixes that chamber and is the identity. Apply this to the coset
geometry to obtain the stated full automorphism group. \(\square\)

\textbf{Lemma C.Gallery (the marked stabilizers).} In a thin regular
residually connected geometry, the full group acts freely on chambers.
At a base chamber let \(r_i\) be the unique group element carrying it
across its i-panel. Then \(r_i\) is an involution, the \(r_i\) generate
the full group, and the stabilizer of a base subflag F is \[
G_F=\langle r_i:i\notin t(F)\rangle.
\tag{C2}
\]

\emph{Proof.} UC.Geometry proves chamber-graph connectedness and
freeness of the chamber action. The square of \(r_i\) fixes the base
chamber, so \(r_i\) has order two. A chamber gallery expresses every
chamber image by successive panel changes, which gives generation. Apply
the same gallery argument in the residue at F. Only the missing types
can change, so its chamber-transitive stabilizer is generated by their
\(r_i\). Conversely those \(r_i\) fix F. This proves (C2). \(\square\)

\textbf{Lemma UC.Helly (anchored criterion).} A finite indexed subgroup
family \((P_i)_{i\in I}\) has typed coset Helly if and only if the
following condition holds. For every \(i\in I\), every
\(J\subseteq I\setminus\{i\}\) with \(|J|\geq2\), and any selected
\(j_0\in J\), each \(g\in P_{j_0}\) satisfying \(g\in P_iP_j\) for all
\(j\in J\) admits \(k\in\bigcap_{j\in J}P_j\) with \(P_i g=P_i k\).

\emph{Proof.} The cosets \(P_i g\) and \(P_j\) intersect precisely when
\(g\in P_iP_j\). Thus Helly gives a common element \(k\), proving
necessity. Conversely, induct on the number of cosets in a pairwise
intersecting family. A family of at most two already has a common
element. For a larger family, the induction hypothesis gives a common
element of all but one coset. Right multiplication by its inverse turns
those cosets into the base cosets \(P_j\). The remaining coset meets
\(P_{j_0}\), so it has a representative \(g\in P_{j_0}\). Its
intersections with the other base cosets are precisely the product
conditions in the statement. The criterion supplies \(k\), and undoing
the right translation gives a common element of the original family.
\(\square\)

\textbf{Lemma C.HellyCount (finite recognition).} For a finite group,
put \[
X_{ij}=\{P_i x:x\in P_j\},\qquad P_J=\bigcap_{j\in J}P_j.
\] Typed coset Helly is equivalent to \[
\left|\bigcap_{j\in J}X_{ij}\right|=[P_J:P_J\cap P_i]
\quad(i\notin J,\ |J|\ge2).
\tag{C3}
\]

\emph{Proof.} The cosets represented in \(P_J\) form a subset of the
displayed intersection, of size the right side. Equality is therefore
equality of these sets. This is exactly the anchored condition of
UC.Helly: a coset meeting every base \(P_j\) has a representative in
their common subgroup. The finite count proves both directions.
\(\square\)

\textbf{Theorem UC.Action (faithful action after retaining objects).}
Let \(G\) act on sets \(X\) and \(Y\). Suppose a \(G\)-invariant subset
\(Y_0\subseteq Y\) has an equivariant surjection onto \(X\). Suppose
also that \(Y\) has an injective equivariant encoding in a tagged
disjoint union of finite products, each factor of which is \(X\) or its
power set, and that the tags are fixed by \(G\). Then \[
\ker(G\curvearrowright Y)=\ker(G\curvearrowright X)
=\bigcap_{x\in X}G_x.
\tag{5.1}
\] If \(X\) is a disjoint union of transitive orbits with selected
stabilizers \(H_j\), the common kernel is \[
N=\bigcap_j\operatorname{core}_G(H_j)
=\operatorname{core}_G\left(\bigcap_jH_j\right).
\tag{5.2}
\] The faithful permutation image is \(G/N\). If its action on \(Y\) is
chamber-transitive on a thin, residually connected finite-rank geometry,
then \(G/N\) is the full type-preserving automorphism group. The
preimages in \(G\) of a chamber stabilizer and of a panel stabilizer are
respectively \(N\) and a subgroup of index two over \(N\).

\emph{Proof.} An element fixing \(Y\) pointwise fixes \(X\) through the
equivariant surjection. Conversely, an element fixing \(X\) fixes all
subsets of \(X\), every finite tuple formed from these objects, and all
fixed tags. Injectivity of the encoding then forces it to fix \(Y\)
pointwise. This proves (5.1). In a transitive orbit all point
stabilizers are conjugates of the selected one, so intersection over all
points gives the first equality in (5.2). The second is the intersection
of the same conjugate subgroups in the opposite order. The first
isomorphism theorem identifies the faithful image.

By UC.Geometry that image has trivial chamber stabilizers and is the
full type-preserving group. Its panel stabilizer acts transitively on
the two chambers through the panel, with trivial stabilizer of either
chamber, so has order two. Taking preimages proves the final assertion.
\(\square\)

Independently labeled copies of the same subset need extra data in their
encoding; they cannot be silently treated as the subset itself.

\textbf{Proposition UC.ExtensionOrbits (classification above a retained
base).} Let a group \(B\) act on sets \(Z,E\), and let \(p:Z\to E\) be
equivariant. For one representative \(e\) of each orbit in
\(B\backslash E\), there is a natural bijection \[
B\backslash Z\ \cong\ \coprod_{e\in B\backslash E}B_e\backslash p^{-1}(e).
\tag{5.6}
\]

\emph{Proof.} Translate any \(z\in Z\) so that \(p(z)\) is the selected
representative of its orbit. If two translated points over that
representative are \(B\)-equivalent, an element taking one to the other
fixes their image \(e\), so belongs to \(B_e\). Conversely a
\(B_e\)-equivalence is already a \(B\)-equivalence. These two
observations prove surjectivity and injectivity of (5.6). \(\square\)

\section{Endpoint conjugation and central
twists}\label{endpoint-conjugation-and-central-twists}

Let \((G;a,b,(r_i)_{i\in E})\) be a C-group marking with E finite and \([a,r_i]=1\)
for all \(i\in E\). No string diagram, finite group, or underlying
geometry is assumed. Define

\[
a'=aba,\quad R_J=\langle r_j:j\in J\rangle,\quad
U_J=\langle a,b,R_J\rangle,\quad H_J=\langle a',b,R_J\rangle,
\] \[
P=\langle b,R_E\rangle,\quad H=H_E,\quad
\mathcal F=\{J\subseteq E:a\in H_J\}.
\tag{E1}
\]

Conjugation by a exchanges b and \(a'\) and fixes \(R_J\) pointwise.
Consequently

\[
H_J\trianglelefteq U_J=H_J\langle a\rangle,
\qquad [U_J:H_J]\le2.
\tag{E2}
\]

\textbf{Theorem C.Endpoint.} The new marking \((H;a',b,(r_i))\) has IP
if and only if

\[
P\cap P^a=R_E,
\qquad \mathcal F\text{ is closed under intersections}.
\tag{E3}
\]

The group H is normal of index at most two in G. It is G precisely when
\(\mathcal F\) is nonempty. Since E is finite, a nonempty
intersection-closed \(\mathcal F\) is precisely the family of supersets
of one least support \(S\subseteq E\).

\textbf{Proof.} The family \(\mathcal F\) is upward closed. Any new
marked subgroup is of one of the forms

\[
A_{\epsilon,J}=\langle b^\epsilon,R_J\rangle\ (\epsilon=0,1),
\qquad H_J,\qquad Q_J^a\text{ with }Q_J=\langle b,R_J\rangle,
\tag{E4}
\]

where \(b^{0}\) means that b is omitted. These correspond to omitting
\(a'\), using both endpoint labels, or using \(a'\) but omitting b.
Intersections of two A's are old marked intersections. Also

\[
A_{\epsilon,I}\cap H_J
\le A_{\epsilon,I}\cap U_J=A_{\epsilon,I\cap J};
\]

all generators on the right belong to \(H_J\), so equality holds.
Intersections of two \(Q^a\)'s follow by conjugating old IP. Conjugation
by a normalizes \(H_I\), so an H-with-\(Q^a\) intersection reduces to
the just proved H-with-A case. For an A-with-\(Q^a\) intersection, the
first condition in (E3) confines it to \(R_E\). Intersecting there and
using that a centralizes \(R_E\) gives \(R_{I\cap J}\).

It remains to intersect \(H_I\) and \(H_J\). Put \(K=I\cap J\). Old IP
gives

\[
H_K\le H_I\cap H_J\le U_I\cap U_J=U_K.
\tag{E5}
\]

If \(a\in H_K\) then \(H_K=U_K\). Otherwise \(H_K\) has index two in
\(U_K\), so the middle subgroup is either \(H_K\) or \(U_K\). The latter
case occurs exactly when a belongs to both \(H_I\) and \(H_J\), and
intersection closure excludes it. This proves all new marked
intersections. Necessity of the first condition follows by intersecting
the two endpoint-deleted new parabolics; necessity of the second follows
by intersecting \(H_I,H_J\). The generators are involutions and the
proved IP supplies independence. Equation (E2) with \(J=E\) gives the
group and index assertions. A finite nonempty upward intersection-closed
family contains its total intersection S, giving the least-support
description. \(\square\)

\textbf{Corollary C.Odd (odd endpoint and proper halving).} If
\(m=|ab|\) is odd, the only test in (E3) is the overlap
\(P\cap P^a=R_E\), and \(H=G\). If H is proper in G, the only test is
again the overlap, and \([G:H]=2\). In either case the endpoint product
has order m/gcd(m,2) for finite m; infinite m remains infinite.

Indeed \(a'b=(ab)^{2}\), and for odd m,

\[
a=(a'b)^{(m+1)/2}b\in H_\varnothing.
\]

Thus \(\mathcal F\) is respectively all subsets or empty. An even
endpoint order by itself does not imply that H is proper; the
least-support alternative in Theorem C.Endpoint cannot be discarded.

\textbf{Lemma C.Twist (central twist).} Let a generating family of a group G contain
involutions b,c with \(|bc|=m\) odd and \(m\ge 3\). Suppose
\(w\in G\) is a nonidentity involution centralizing b and c. Replacing
b by wb in that family preserves G and changes \(|bc|\) to \(2m\).

\textbf{Proof.} The commuting factors w and bc have coprime orders, so
\((wbc)^m=w\) and \(|wbc|=2m\). Hence the new generators recover w and
b. The new marker wb is a nonidentity involution: \(wb=1\) would make
\(b=w\) commute with c, contrary to \(m\ge 3\) odd. \(\square\)

A parabolic containing w,b,c is unchanged by this twist; one omitting b
is unchanged. Other intersections and typed Helly still require proof.
Lemma C.Twist is the auxiliary operation needed in the finite
application below, not a blanket geometric preservation theorem.

\section{Graph halving and lattice
polytopes}\label{graph-halving-and-lattice-polytopes}

Let \(\Gamma\) be a finite-rank residually connected incidence geometry,
of rank \(n\ge 3\). Designate types 0 and 1 as vertices and edges.
Assume their truncation is a simple graph Q, and that an edge is
incident with a higher object exactly when both its endpoints are
incident with that object. For a flag F of higher objects, let \(Q_F\)
be its common vertex-edge shadow. All \(Q_F\) are connected and contain
an edge.

UC.Geometry supplies the connected chamber graph of every residually
connected finite-rank geometry. Restricting a chamber gallery to any set
of at least two types shows that the corresponding truncation is
connected, because successive complete incidence subgraphs overlap.
Apply this inside the residue at F; chamber extension supplies an edge.
The edge-shadow hypothesis says \(Q_F\) is the induced graph on its
common vertices.

For a connected graph Q, its bipartite double cover D(Q) has vertices
\((v,\varepsilon )\), \(\varepsilon \in \{0,1\}\), with
\((v,\varepsilon )\) adjacent to \((w,1-\varepsilon )\) when v is
adjacent to w. It is connected when Q has an odd cycle, and has two
components otherwise: parity of walks proves both assertions. Choose a
component C of D(Q). Define \(L_C(\Gamma )\) as follows. Its first two
types are (v,0) and (v,1) in C. A higher object is a pair \((x,C_x)\),
where \(C_x\) is a component of \(D(Q_x)\) contained in C. Incidence of
the first two types is adjacency, incidence with a higher object is
membership in \(C_x\), and two higher objects are incident when their
old objects are incident and their chosen components meet.

The precise compatibility requirement is

\[
\text{if every }Q_x\ (x\in F)\text{ is nonbipartite, then }Q_F
\text{ is nonbipartite},
\tag{PC}
\]

for every nonempty higher flag F. It is automatic if the higher shadows
of every flag are nested, or if Q is bipartite.

\textbf{Theorem C.Graph.} The system \(L_C(\Gamma )\) is always a
geometry. It is residually connected exactly when (PC) holds. It is thin
exactly when \(\Gamma\) is thin. A chamber-transitive automorphism group
of \(\Gamma\) acts faithfully and chamber-transitively on
\(L_C(\Gamma )\) when Q is nonbipartite; when Q is bipartite its
index-two color-preserving subgroup does so.

\textbf{Proof.} For a flag \(\Phi =\{(x,C_x):x\in F\}\) of higher lifts,
write \(W=\cap C_x\). If any \(Q_x\) is bipartite, restriction of
\(C_x\) to \(D(Q_F)\) is one component. Any two such restrictions agree:
incidence makes them meet over the connected \(Q_{x,y}\), and hence
select the same parity class there and over \(Q_F\). Nonbipartite
constituents impose no restriction. Thus W is one component if any
constituent is bipartite, and all of \(D(Q_F)\) otherwise. In particular
W always contains an edge, and is connected exactly under (PC) for this
flag. For F empty use \(W=C\).

An edge of W and a chamber of \(\Gamma\) extending its base edge and F
give a chamber extending \(\Phi\). Every additional higher object has a
unique component containing the chosen lift of that edge. The same
extension works for a flag already containing one lifted vertex. If it
contains both, their base edge is unique by simplicity and incident with
all the old higher objects by the shadow condition. This proves the
geometry axiom.

The residue at (v,0) is the old residue \(\Gamma _v\): a neighbour (w,1)
corresponds to the unique edge \{v,w\}, and the component for an
incident higher object is uniquely specified by (v,0). Simplicity and
the shadow condition preserve all incidences. The other color is
identical. Hence residues of flags containing a lifted vertex are
connected. For a flag of higher objects only, its two-type truncation is
W; when W is connected, every other residual object attaches to it
through a chamber. This proves residual connectivity under (PC). If (PC)
fails, choose its higher flag; all constituents have unique lifts but W
has two components. Were the resulting geometry residually connected,
the truncation lemma inside its residue would make W connected. This
proves necessity without treating a truncation as a residue.

Panels containing both new vertex types correspond to old panels missing
the same higher type. A panel missing one new vertex type corresponds to
an old panel missing the edge type, at the retained endpoint. Old panels
missing the vertex type already have two objects by simplicity. These
correspondences prove both directions of thinness.

If Q is nonbipartite, an old chamber
\((v,\{v,w\},x_{2},\ldots ,x_{n-1})\) corresponds bijectively to the new
chamber with (v,0),(w,1) and the components containing them. Old
automorphisms act naturally on this bijection. If Q is bipartite, the
same bijection uses the old chambers whose chosen vertex is in one fixed
color class. A chamber-transitive group contains a color-exchanging
element, since it exchanges the endpoints of a base edge through a
chamber move. Its color-preserving subgroup therefore has index two and
is transitive on the selected chambers, with the same chamber bijection.
Faithfulness in both cases follows from the retained-object argument
below. \(\square\)

\subsection{The retained objects}\label{the-retained-objects}

For completeness, apply UC.Action to the precise object sets in this
construction. Let K be G when Q is nonbipartite and its color-preserving
subgroup otherwise. Take X to be the disjoint union of old vertices and
higher objects, and Y the object set of the halved geometry. The map \[
(v,\varepsilon)\mapsto v,\qquad (x,C_x)\mapsto x
\tag{C4}
\] is an equivariant surjection \(Y\to X\). For nonbipartite Q the
chosen C is all D(Q); for bipartite Q it contains exactly one lift of
every old vertex. In the latter case its intersection with the double
cover of each connected shadow \(Q_x\) is one of that double cover's two
components. Thus every higher object also has a lift.

Encode a new vertex by its fixed color tag and old vertex. Encode
\((x,C_x)\) by the old object x and the two subsets
\(\{v:(v,0)\in C_x\}\), \(\{v:(v,1)\in C_x\}\) of X, retaining its old
type as a tag. Components are determined by their vertex sets, so the
encoding is injective and equivariant. UC.Action gives \[
\ker(K\curvearrowright Y)=\ker(K\curvearrowright X)=1.
\tag{C5}
\] The last equality follows because an old edge is uniquely determined
by its two endpoints in the simple graph Q. Fixing X therefore fixes all
old objects. When \(\Gamma\) is thin and (PC) holds, the proved geometry
is thin and residually connected, and its K-action is
chamber-transitive. UC.Geometry then identifies K with the full
type-preserving group. This argument does not infer the geometry or its
flag properties merely from faithfulness.

Suppose now \(\Gamma\) is regular and thin, with vertex-changing marker
a, edge-changing marker b, and higher markers \(r_i\); assume a commutes
with the higher markers. Under (PC), Theorem C.Graph produces a regular
thin residually connected geometry. Directly on a base edge, its panel
involutions are \(b^a,b,r_i\): b changes one endpoint, aba changes the
other, and each \(r_i\) changes its higher object. UC.Geometry
identifies their group as the full type-preserving group H. Lemma
C.Gallery identifies the stabilizer of every base subflag with the
subgroup generated by its missing-type markers, giving IP and
identifying the actual coset incidence with \(L_C(\Gamma )\).

There is also an exact correspondence between the two parity conditions.
In the base residue with group \(U_J\) from (E1), the vertex stabilizer
is \(Q_J=\langle b,R_J\rangle\). Its vertex graph is bipartite precisely
when \(a\notin H_J\). If \(H_J\) is proper, its two cosets give a
coloring, since \(Q_J\le H_J\) and a crosses colors. Conversely the
action on the two colors supplies a character killing b and all \(r_j\)
and sending a to the nontrivial element, so it excludes a from \(H_J\).
In a globally nonbipartite geometry, higher flags reduce by transitivity
to base flags with shadows \(U_{E\setminus S}\). Condition (PC) says

\[
[E\setminus\{i\}\in\mathcal F\text{ for all }i\in S]
\Longrightarrow E\setminus S\in\mathcal F.
\tag{G1}
\]

An upward family containing E satisfies (G1) exactly when it is
principal: let T be the set of i whose complementary coatom belongs to
\(\mathcal F\); (G1) gives \(E\setminus T\in \mathcal F\), and every
member of \(\mathcal F\) contains \(E\setminus T\). Thus geometric parity
compatibility is the least-support condition of Theorem C.Endpoint. The
mixed overlap in that theorem has a different role: it says the two
adjacent vertices determine one edge. Indeed P and \(P^a\) stabilize the
endpoints; their intersection is the edge's pointwise stabilizer \(R_E\)
exactly when no other edge has those endpoints, using transitivity on
panels and the regular chamber identification.

\subsection{Halving lattice polytopes}\label{halving-lattice-polytopes}

\textbf{Corollary C.Halving.} Let P be a regular rank-n abstract
polytope, \(n\ge 3\), whose face poset is a lattice. Replacing one
endpoint marker a by aba gives a C-group H normal of index at most two
in Aut(P). Its maximal-parabolic coset system is a regular hypertope,
with full type-preserving group H. The diagram duplicates the edge
labels from b to the remaining markers, and the \(a'–b\) label is
m/gcd(m,2) for \(m=|ab|\) finite, or infinity for infinite m. The index
is two exactly when the corresponding vertex graph is bipartite; for a
final-end operation use the dual graph on facets.

\textbf{Proof.} In a lattice polytope an edge is the join of its two
distinct vertices. Thus two vertices determine at most one edge, and
that edge belongs to a higher face exactly when both vertices do. Higher
faces of a flag form a chain, so their shadows are nested and (PC) is
automatic. Theorem C.Graph and the panel-marker identification prove the
asserted geometry, full group, and IP without an
alternating-semiregular-polytope theorem. Equation (E2) and the coloring
argument give the index. Finally \(a'b=(ab)^{2}\) and a commutes with
all \(r_i\), so \(a'r_i\) is conjugate to \(br_i\); all other products
are unchanged. \(\square\)


\section{Graph differences and even subsets}
For a nonempty set $\Omega$, let $\F_2^{(\Omega)}$ be the vector space
of finitely supported functions on $\Omega$, with coordinate vectors
$e_v$. Its augmentation module is
\[
 E(\Omega)=\left\{z\in\F_2^{(\Omega)}:
                 \sum_{v\in\Omega}z_v=0\right\}.
\]
For a chosen vertex $o$, the vectors $e_o+e_v$, $v\ne o$, are a basis.
Indeed a vector of even total sum is the sum of $e_o+e_v$ over its
nonzero coordinates away from $o$; independence follows by inspecting
those coordinates. Thus, if $|\Omega|=m<\infty$, the module has order
$2^{m-1}$.

\begin{theorem}[Presentation by graph differences]\label{UC.graphmodule}
Let $Q$ be a connected simple graph with vertex set $\Omega$. Define
$B(Q)$ by one generator $z_e$ for every edge, with relations
\[
 z_e^2=1,\qquad [z_e,z_f]=1,
 \qquad z_{e_1}\cdots z_{e_k}=1
\]
for every finite closed walk with successive edges $e_1,\ldots,e_k$.
Then $z_{\{u,v\}}\mapsto e_u+e_v$ defines an isomorphism
$B(Q)\cong E(\Omega)$. It commutes with every graph automorphism.
\end{theorem}
\begin{proof}
The edge vectors have order dividing two and commute. Their sum along a
closed walk is zero because each successive vertex occurs twice. Hence
the indicated homomorphism exists.

Choose a base vertex $o$. For each $v$, choose a finite path from $o$ to
$v$ and let $t_v$ be the product of its edge generators. The product is
independent of the chosen path: concatenate one path with the reverse
of the other, and apply the closed-walk relation. In particular $t_o=1$.
Following a path to $u$, then the edge $\{u,v\}$, and finally a reverse
path from $v$ to $o$ gives $z_{\{u,v\}}=t_ut_v$.

The elements $t_v$ commute and have square one. The basis description
above therefore defines a homomorphism $E(\Omega)\to B(Q)$ sending
$e_o+e_v$ to $t_v$. The two composites fix every edge generator and
every basis vector, respectively, so they are identity maps.
The original edge map, unlike the chosen inverse formula, uses no
base vertex or paths. It consequently commutes with every graph
automorphism.
\end{proof}

For a finite graph one may use a spanning tree and only the fundamental
cycles determined by the remaining edges. The same proof applies:
the unique tree paths define $t_v$, and each fundamental-cycle relation
gives $z_{\{u,v\}}=t_ut_v$ for its extra edge. This also proves directly
that every closed-walk relation follows from that finite list.

\begin{corollary}[Generation inside a permutation extension]\label{UC.edgegeneration}
Let $H$ act by graph automorphisms on a connected graph $Q$. In
$E(\Omega)\rtimes H$, suppose a subgroup contains $\{0\}\rtimes H$
and translations $e_u+e_v$ for a set of edges whose $H$-orbits form a
connected spanning subgraph. Then it is the whole semidirect product.
\end{corollary}
\begin{proof}
Conjugation by $(0,h)$ carries the translation $e_u+e_v$ to
$e_{hu}+e_{hv}$. The subgroup therefore contains every edge difference
in the indicated spanning subgraph. Theorem~\ref{UC.graphmodule} says
that these generate $E(\Omega)$, proving the assertion.
\end{proof}

\section{Recovering a semidirect product from its translations}
Let $H$ act on a nonempty set $\Omega$, and fix $o\in\Omega$.
We use left actions and multiplication
\[
 (v,g)(w,h)=(v+gw,gh)
       \quad\hbox{in }E(\Omega)\rtimes H.
\]
The following statement does not require either set or group to be
finite.

\begin{theorem}[Recognition by differences from a base vertex]\label{UC.recognition}
Let $P$ be a group, let $j:H\to P$ be a homomorphism, and let
$(t_v)_{v\in\Omega}$ be commuting elements of $P$ satisfying
$t_v^2=1$ and $t_o=1$. Suppose
\begin{equation}\label{UC.covariance}
 j(h)t_vj(h)^{-1}=t_{ho}t_{hv}
        \qquad(h\in H,\ v\in\Omega).
\end{equation}
There is a unique homomorphism
\[
 \Psi:E(\Omega)\rtimes H\longrightarrow P,
 \qquad \Psi(e_o+e_v,1)=t_v,
 \qquad \Psi(0,h)=j(h).
\]
It is surjective if these elements generate $P$. Under that generation
hypothesis, if a homomorphism $\Phi:P\to E(\Omega)\rtimes H$ has
the displayed inverse values on these elements, then $\Phi$ and
$\Psi$ are mutually inverse.
\end{theorem}
\begin{proof}
The commuting involutions define a homomorphism
$T:E(\Omega)\to P$ by $T(e_o+e_v)=t_v$ on the stated basis.
The action of $h$ on this vector is
\[
 h(e_o+e_v)=(e_o+e_{ho})+(e_o+e_{hv}).
\]
Consequently~\eqref{UC.covariance} implies
$T(hw)=j(h)T(w)j(h)^{-1}$ first on every basis vector and then on
every vector. Define $\Psi(w,h)=T(w)j(h)$. For two such pairs,
\[
 T(v)j(g)T(w)j(h)=T(v)T(gw)j(gh)=T(v+gw)j(gh),
\]
which is exactly the required homomorphism identity. Uniqueness and
surjectivity follow from the generating sets. Under the final
hypothesis each composite fixes its respective generating set, and is
therefore the identity.
\end{proof}

\begin{lemma}[A twist exchanging two complements]\label{UC.twist}
The function $c(h)=e_o+he_o$ satisfies
$c(gh)=c(g)+gc(h)$. Thus
\[
 \delta(v,h)=(v+c(h),h)
\]
is an involutory automorphism of $E(\Omega)\rtimes H$. The standard
complement $A=\{(0,h):h\in H\}$ and its image satisfy
\[
 A\cap\delta(A)=\{(0,h):h\in H_o\},
\]
where $H_o$ is the stabilizer of $o$.
\end{lemma}
\begin{proof}
In characteristic two,
$e_o+ghe_o=(e_o+ge_o)+g(e_o+he_o)$, proving the cocycle identity.
Substitution into the multiplication rule proves that $\delta$ is a
homomorphism; applying it twice adds $c(h)$ twice, so gives the identity.
Finally $(0,h)$ belongs to its image of $A$ exactly when
$c(h)=0$, or equivalently $ho=o$.
\end{proof}

\section{Cyclic displacement relations}
The next identity is the displacement argument in Li--Zhang's proof of
their Theorem~3.4. Keeping integer indices exposes the action needed in
Theorem~\ref{UC.recognition} and avoids a separate order argument.

\begin{lemma}\label{UC.displacements}
Let $x,y$ belong to a group and put $t_k=y^kx^k$ for $k\in\Z$.
If every $t_k$ has square one, then all $t_k$ commute and
\[
 x^{-1}t_kx=t_1t_{k+1},\qquad
 xt_kx^{-1}=t_{-1}t_{k-1}.
\]
If an involution $b$ inverts $x$ and $y$, then $bt_kb=t_{-k}$.
If $x^n=y^n=1$, it suffices to require $t_k^2=1$ for
$1\le k<n$.
\end{lemma}
\begin{proof}
From $t_k=t_k^{-1}$ we obtain
$y^kx^k=x^{-k}y^{-k}$ and, on conjugating by $x^k$,
$x^ky^k=y^{-k}x^{-k}$. Hence
\[
 \begin{aligned}
 t_it_j
 &=x^{-i}y^{j-i}x^j
 =x^{-j}(x^{j-i}y^{j-i})x^j\\
 &=x^{-j}y^{i-j}x^i=t_jt_i.
 \end{aligned}
\]
Direct multiplication gives
$x^{-1}t_kx=t_1^{-1}t_{k+1}=t_1t_{k+1}$.
For $k=-1$ this yields $xt_1x^{-1}=t_{-1}$.
Conjugate the first formula with index $k-1$ by $x$ and use this
identity to obtain the second formula. Conjugation by $b$ gives
$bt_kb=y^{-k}x^{-k}=t_{-k}$. Finally the two period relations give
$t_{k+n}=t_k$ for every integer $k$, while $t_0=1$; this proves the
finite-relation assertion.
\end{proof}

\begin{corollary}\label{UC.cyclicrecognition}
Suppose $n\ge3$ and $P=\gen{x,b,y}$ satisfies
\[
 x^n=y^n=b^2=1,\quad bxb=x^{-1},\quad byb=y^{-1},
 \quad(y^kx^k)^2=1\quad(1\le k<n).
\]
Let $H=D_{2n}$ act on $\Z/n\Z$ by $a(i)=i+1$ and $b(i)=-i$.
If there is a homomorphism $\Phi:P\to E(\Z/n\Z)\rtimes H$ with
\[
 x\longmapsto(0,a),\quad b\longmapsto(0,b),\quad
 y\longmapsto(e_0+e_{-1},a^{-1}),
\]
then it is an isomorphism. Its inverse sends
$(e_0+e_i,1)$ to $t_{-i}$ and $(0,a),(0,b)$ to $x,b$.
\end{corollary}
\begin{proof}
The dihedral relations define $j:H\to P$, with $j(a)=x$ and
$j(b)=b$. Set $T_i=t_{-i}$; the period relations make this independent
of the integer representative of $i$. Lemma~\ref{UC.displacements}
gives commuting involutions and
\[
 xT_ix^{-1}=T_1T_{i+1},\qquad bT_ib=T_{-i}.
\]
These are exactly~\eqref{UC.covariance} for $a,b$, with $o=0$.
Equivariance for products and inverses follows by applying these
identities successively; thus it holds for all $h\in H$.
The elements $j(H),T_i$ generate $P$, since $y=t_1x^{-1}$.
Moreover $\Phi(t_k)=(e_0+e_{-k},1)$ by direct multiplication, or by
the cocycle identity in Lemma~\ref{UC.twist}. The final assertion of
Theorem~\ref{UC.recognition} now gives the inverse and the isomorphism.
\end{proof}

\section{Formal components and their limits}
The file \texttt{SemidirectRecognition.lean} uses mathlib's standard
semidirect-product construction. It verifies the actual lift
$\Psi(n,h)=T(n)j(h)$, proves its surjectivity from the stated generation
hypothesis, and constructs the multiplicative equivalence from the two
component identities. Its nine reported declarations make no
finiteness assumption. The augmentation basis and the passage from
cyclic displacement generators to these component homomorphisms are
written arguments above, rather than implicit formal hypotheses.

Two unchanged files are also included: the five reported displacement
declarations in \texttt{BinaryDihedralDisplacements.lean}, and the two
mixed-intersection declarations in \texttt{EndpointDoublingIntersections.lean}.
The latter assumes the fixed-overlap hypotheses of the appropriate
case of C.Endpoint. These components do not formalize the complete
endpoint criterion, coset geometry, graph halving, or source-specific
applications. The separate checker recompiles all three files and
reports these limits explicitly.
\section*{References}
{[}MW20{]} A. Montero and A. Ivi\'c Weiss,
\emph{Locally spherical hypertopes from generalised cubes},
arXiv:2001.11082v1, 29 January 2020.

{[}LZ25{]} M. Li and W.-J. Zhang,
\emph{Regular 3-polytopes of type $\{n,n\}$},
arXiv:2505.09505v1, 14 May 2025.

{[}PT24{]} C. A. Piedade and P. Tranchida,
\emph{Constructing new geometries: a generalized approach to halving
for hypertopes}, arXiv:2405.19050v1. The graph construction and lattice
convention are used at the explicit scope stated above.
\end{document}
